% ============================================================
%  Chapter 01: The Sociological Perspective, Modernity, 
%  Social Changes in Europe & Emergence of Sociology
% ============================================================

\chapter{The Sociological Perspective, Modernity \& Social Changes in Europe}

\section{Developing the Sociological Perspective}

\subsection{Impression Management}
The instructor introduced the idea that every social role is a performance . He said that for the duration of the course, both teacher and students would be ``acting'' their roles---he as teacher, they as students---and that each side is cautious in how it behaves before the other in order to maintain this arrangement .

\begin{KeyConcept}
\begin{itemize}
    \item This behaviour of consciously managing how one appears to others in a given role is called \textbf{impression management} . Scholar cited: \textbf{Erving Goffman} .
    \item The instructor illustrated it with his own multiple roles: he is a teacher to his students, but a student to his own PhD supervisor, and a son to his mother---no single role captures him ``in absolute sense'' .
\end{itemize}
\end{KeyConcept}

\begin{QuickRevision}
\begin{itemize}
    \item \textbf{Impression management (Goffman):} consciously managing how one appears in a social role; no one occupies a single role ``in absolute sense'' .
\end{itemize}
\end{QuickRevision}

\subsection{Seeing the ``System'', Not the Individual}
The instructor ran a picture-based exercise to train students to look for the underlying system, force, or structure behind what is visible, rather than describing only the visible individuals or objects . This is presented as the foundational move of sociological thinking .

\begin{ExampleBox}
\begin{itemize}
    \item \textbf{Picture 1 --- The Solar System:} Students named the planets they saw, but the instructor pointed out that what is really being seen is a system---the sun giving light to other planets---not any individual planet .
    \item \textbf{Picture 2 --- A Man Throwing/Holding a Ball:} Students gave varied answers (a man, a ball, projectile motion), but the deeper common thread the instructor drew out was gravitational force / centripetal force---something that keeps bodies in motion but is not itself empirically visible; it can only be felt, not seen, touched or heard .
    \item \textbf{Picture 3 --- A Benzene Ring (Six Carbons, Six Hydrogens):} As a student of chemistry one does not describe six carbons and six hydrogens individually but names the structure as a whole---benzene---which is something different from the sum of the individual atoms .
\end{itemize}
\end{ExampleBox}

\begin{CriticalPoint}
\begin{itemize}
    \item In each case the instructor did not talk about individual planets, individual men, or individual atoms---only about the system, force, or structure . He said this same shift (from individual to structure) is what sociological perspective requires throughout the picture exercise that follows .
\end{itemize}
\end{CriticalPoint}

\begin{QuickRevision}
\begin{itemize}
    \item \textbf{Sociological perspective:} Seeing the system/force/structure (not the individual)---illustrated with the solar system, gravity, and benzene .
\end{itemize}
\end{QuickRevision}

\subsection{Conspicuous Consumption}
Shown a picture of a Starbucks coffee cup, students named the brand, the cup, the price . The instructor's sociological reading:

\begin{KeyConcept}
\textbf{Conspicuous consumption}---buying an expensive good not for its use-value but to display one's status .
\begin{itemize}
    \item Economics predicts that demand falls as price rises (the law of demand), yet Starbucks coffee, though costly, is in high demand---people prefer it over cheaper options such as Indian Coffee House .
    \item The instructor's explanation: the buyer is not purchasing coffee but recognition---the moment of swiping a card, or taking a selfie with the cup, displays one's class status to others .
\end{itemize}
This is presented as a challenge to the economic idea of man as \textit{homo economicus} (a purely utilitarian, rational being)---in conspicuous consumption, the utility being sought lies beyond the commodity itself, in the status it signals .
\end{KeyConcept}

\begin{QuickRevision}
\begin{itemize}
    \item \textbf{Conspicuous consumption:} Buying to display status, not for use-value; challenges the idea of man as purely rational \textit{homo economicus} (Starbucks example) .
\end{itemize}
\end{QuickRevision}

\subsection{Poverty and Structural Inequality}
Shown a picture of a mother and child in poor conditions, students' answers ranged from ``poverty'' and ``happy despite hardship'' to ``malnutrition'' and ``single motherhood'' .

\begin{KeyConcept}
\begin{itemize}
    \item As a sociologist, the instructor said, one does not see ``a woman'' and ``a child'' but inequality, division of society into economic classes, unequal opportunity, unemployment, single motherhood, and malnutrition---conditions read as structural rather than individual .
\end{itemize}
\end{KeyConcept}

\begin{QuickRevision}
\begin{itemize}
    \item \textbf{Poverty read structurally:} Inequality, unequal opportunity, unemployment, single motherhood, malnutrition---not individual failing .
\end{itemize}
\end{QuickRevision}

\subsection{Child Labour and Cheap Labour}
A picture of children engaged in labour drew answers such as ``child labour,'' ``exploitation,'' ``alienation'' and ``absence of human rights'' .

\begin{CriticalPoint}
\begin{itemize}
    \item The instructor added a second, structural reading: from the position of an industrialist, the same picture shows the cheapest labour available---the same phenomenon looks different depending on the social location (industrialist vs sociologist) from which it is read .
\end{itemize}
\end{CriticalPoint}

\begin{QuickRevision}
\begin{itemize}
    \item \textbf{Child labour:} Exploitation/alienation to a sociologist, but reads as ``cheapest labour available'' from an industrialist's social location---same fact, different vantage point .
\end{itemize}
\end{QuickRevision}

\subsection{Division of Labour and Sexual Division of Labour}
Shown a picture of a man working on a laptop and a woman washing dishes, students named ``work from home,'' ``patriarchal society,'' and ``sexism'' . The instructor built up the underlying concept step by step .

\begin{KeyConcept}
\textbf{Division of labour}, illustrated through McDonald's: a burger is ready in about two minutes because five separate people specialise in five separate tasks (making the tikki, dressing the bun, frying the fries, refilling the coke, etc.)---specialised roles integrated together make production efficient .

\textbf{Scholar and text:} Adam Smith, \textit{Wealth of Nations}---an inquiry into why Britain became prosperous . Smith's answer: division of labour . Industrial revolution brought technological advancement and mass production, but mass production itself was possible only because of division of labour .
\end{KeyConcept}

\begin{KeyConcept}
\textbf{Sexual division of labour}---when tasks are assigned not on the basis of skill but on the basis of gender . A woman is as capable of being trained to work on a laptop as a man; the reason she is instead confined to housework is not lack of ability but social assignment of roles by sex .
\end{KeyConcept}

\begin{QuickRevision}
\begin{itemize}
    \item \textbf{Division of labour (Adam Smith, \textit{Wealth of Nations}):} Specialised tasks for efficiency .
    \item \textbf{Sexual division of labour:} Tasks assigned by gender, not skill .
\end{itemize}
\end{QuickRevision}

\subsection{Media and the Reproduction of Patriarchy}
A screenshot instructing ``how to be a good housewife'' (cleaning, cooking, cutting vegetables) was used to make a point about media as a system, not as a single website .

\begin{CriticalPoint}
\begin{itemize}
    \item The instructor's reading: media as a whole (not one website) promotes the same gender-stereotypical roles assigned to different sexes; media is described as being ``in alliance with patriarchy,'' each reproducing and reinforcing the other and keeping women confined to prescribed roles .
    \item Just as gravitational force and patriarchy are both not directly visible but can be ``felt'' in their effects, patriarchy is described as an ideology that cannot itself be seen, only its consequences can .
\end{itemize}
\end{CriticalPoint}

\begin{QuickRevision}
\begin{itemize}
    \item Media (as a system, not one website) is in alliance with patriarchy, reproducing gender-stereotypical roles; patriarchy, like gravity, is felt through its effects though not itself directly visible .
\end{itemize}
\end{QuickRevision}

\subsection{Alienation, Role Diversity, and Emotional Labour}
A picture of women doing repetitive garment/mask-stitching work drew initial answers of ``skilling'' and ``women empowerment'' . The instructor complicated this reading .

\begin{CriticalPoint}
\begin{itemize}
    \item The instructor questioned whether monotonous, repetitive factory work over 10 hours a day with no ability to look around, chat, or vary the task can be called empowerment, contrasting it with the role diversity available to, say, an IAS officer .
    \item \textbf{Concept invoked: Alienation (Karl Marx).} With the industrial revolution, man loses control over his productive skills; every day is a repetition of the same task with no colour or charm left in work .
    \item Marx's idea of \textbf{species being} was cited: human beings have wide-ranging potential (e.g., a mother who is nurturing may have many other talents), but the way society and the economic system are organised confines individuals---such as this mother, or the assembly-line worker---to a single narrow role, and that potential is lost .
\end{itemize}
\end{CriticalPoint}

\paragraph{Emotional Labour}
Discussing service-sector roles such as hotel staff and airline cabin crew, the instructor introduced a further concept .

\begin{KeyConcept}
\textbf{Emotional labour}---smiling and being warm to customers even when one does not feel like it, e.g., cabin crew smiling for hours regardless of their actual mood . Smiling without genuine wish to is described as a form of alienation from one's own self .
\begin{itemize}
    \item The instructor offered two readings of the service sector in the same breath: it has expanded opportunities for many groups (via internet-based and remote work), while at the same time emotional labour has made the condition of human existence, in his words, ``more miserable'' . He noted both readings can be right at once---sociology is about seeing things from different angles, not declaring one right and one wrong .
\end{itemize}
\end{KeyConcept}

\begin{QuickRevision}
\begin{itemize}
    \item \textbf{Alienation (Marx):} Specialised, monotonous work strips role diversity and loses man's ``species being'' potential .
    \item \textbf{Emotional labour:} Forced smiling in service work alienates one from one's own self .
\end{itemize}
\end{QuickRevision}

\subsection{Gender, Occupation and the Military}
A picture of army/military personnel drew responses of ``nationalism,'' ``unity,'' ``masculinity'' and ``military strength'' .

\begin{CriticalPoint}
\begin{itemize}
    \item The instructor's reading: \textbf{occupational segregation by sex}---the image ``army officer'' conjures a man in the popular imagination, not a woman, even though bravery and courage are not male-exclusive traits .
    \item \textbf{Exam-relevance note:} The instructor referenced the then-ongoing issue of permanent commission for women in the armed forces as an example of contested access to positions associated with bravery/valour .
    \item He also pointed to healthcare and frontline pandemic workers---Asha workers, Anganwadi workers, auxiliary nurse midwives (ANMs), lady doctors---as the ``real warriors'' of the (then) crisis, most of whom are women, arguing this challenges the assumption that only men can fight for the nation .
\end{itemize}
\end{CriticalPoint}

\begin{QuickRevision}
\begin{itemize}
    \item Occupational imagery is gendered by default (``army officer'' evokes a man); frontline/healthcare women workers reframed as ``real warriors'' .
\end{itemize}
\end{QuickRevision}

\subsection{Voluntary Segregation and Gated Communities}
A picture of a residential gate/quarters drew answers naming class segregation .

\begin{KeyConcept}
The instructor pointed to how residences of IAS/IPS officers, chief ministers, the President's estate, and five-star hotels are all physically excluded from the mainstream city---a spatial way in which status and class segregation are expressed .

\textbf{Term introduced: Gated community}---voluntary exclusion by some groups from mixing with ``ordinary'' people, a further form of inequality since not everyone can afford to belong to one .
\end{KeyConcept}

\begin{QuickRevision}
\begin{itemize}
    \item \textbf{Gated community:} Voluntary spatial self-exclusion by elites (IAS/IPS bungalows, President's estate)---a further form of inequality, since not everyone can afford it .
\end{itemize}
\end{QuickRevision}

\subsection{Schooling, Discipline and the Hidden Curriculum}
A picture of a classroom drew answers such as ``discipline,'' ``socialization,'' and ``value system'' .

\begin{KeyConcept}
\textbf{Hidden curriculum}---the unstated lessons a school teaches beyond its official subjects (maths, physics, English): authority, obedience, conformity, and fear of punishment . A child is taught not to be ``too creative'' and not to question beyond the classroom .
\begin{itemize}
    \item The instructor drew a direct parallel between the classroom and the factory: school prepares different students for the future roles of manager, supervisor, or worker, and in that sense school is described as ``the friend of capitalism''---not only in an economic sense but through the discipline it instils .
    \item Before school, a child's only authority is the parents; school introduces, for the first time, a controlling authority other than the parents---the teacher .
\end{itemize}
\end{KeyConcept}

\begin{QuickRevision}
\begin{itemize}
    \item \textbf{Hidden curriculum:} School teaches authority, obedience and conformity alongside its official subjects; school is described as preparing students for factory-like manager/supervisor/worker roles .
\end{itemize}
\end{QuickRevision}

\subsection{Surveillance and Conformity}
A picture of a CCTV camera drew the answer ``surveillance'' .

\begin{CriticalPoint}
\begin{itemize}
    \item The instructor's elaboration: the function of surveillance cameras is to create fear among those who might otherwise commit crime (theft, snatching, etc.), which minimises criminal tendencies and reproduces conformity and discipline---this is presented as one way society avoids deviance .
    \item He also distinguished \textbf{equality of opportunity} from \textbf{equality of outcome}, noting that opportunity can be made equal, but outcomes will always differ .
\end{itemize}
\end{CriticalPoint}

\begin{QuickRevision}
\begin{itemize}
    \item Surveillance produces fear of detection, minimises deviance, and reproduces conformity; equality of opportunity is distinguished from (unattainable) equality of outcome .
\end{itemize}
\end{QuickRevision}

\subsection{Social Stratification: Caste, Elites, and the Idea of Revolution}
A picture symbolising social hierarchy drew answers of ``social structure'' and ``social hierarchy'' . The instructor used it to introduce the idea of an estate/state system (stratification) and to bridge into the historical material that follows .

\begin{GovtPolicy}
The Indian caste system was cited as an illustration of stratification: Brahmins at the top, Kshatriyas second, Vaishyas third, and Shudras fourth, with a fifth stratum existing beyond these four---opportunity in this system is tied to one's position in the hierarchy .

This was directly compared to pre-revolutionary French society, which the instructor described as similarly stratified, and to the challenge that stratified order eventually faced: the Tennis Court Oath, described here as a declaration that man has an individual existence of his own that religious or hierarchical authority ``in the name of God'' could no longer override .
\end{GovtPolicy}

\begin{PYQLink}
\begin{itemize}
    \item This picture-based comparison of Indian caste stratification with the French estate system is flagged by the instructor as a bridge into the historical unit on Modernity and Social Changes in Europe that follows .
\end{itemize}
\end{PYQLink}

\begin{QuickRevision}
\begin{itemize}
    \item Indian caste stratification ($\text{Brahmin} \rightarrow \text{Kshatriya} \rightarrow \text{Vaishya} \rightarrow \text{Shudra}$) compared to pre-revolutionary French estates; Tennis Court Oath cited as the historical challenge to such hierarchical order---the bridge into Topic 2 .
\end{itemize}
\end{QuickRevision}

\subsection{Family, Heteronormativity and Deviance}
A picture of a conventional family unit drew the answer ``happy family'' . The instructor introduced a further concept .

\begin{KeyConcept}
\textbf{Heteronormativity}---the assumption, embedded in such images, that only heterosexual pairing (man-woman) constitutes a family, so that families formed by same-sex couples are not socially recognised as families, even though assisted/artificial reproductive technologies (ART) now make biological compulsion for children unnecessary .
\end{KeyConcept}

A subsequent picture depicting a same-sex couple/family drew answers of ``homosexuality,'' ``LGBT,'' and ``rebellion'' .

\begin{KeyConcept}
\textbf{Deviance}---breaking away from the standard, socially expected notion (here, of family) . The instructor noted the term carries a negative connotation in common usage (``breaking away''), and that this image can be read as both a challenge to heteronormative standards and, in a value-neutral sociological sense, as deviance from the dominant social pattern .
\end{KeyConcept}

\begin{QuickRevision}
\begin{itemize}
    \item \textbf{Heteronormativity:} Family imagery assumes only heterosexual pairing counts as family, despite ART removing any biological need for this .
    \item \textbf{Deviance:} Breaking from a socially expected pattern---read here as both ``challenge'' and, value-neutrally, deviation from the dominant norm .
\end{itemize}
\end{QuickRevision}

\sectiondivider

\section{Sociology, the Discipline: Modernity, Social Changes in Europe \& Emergence of Sociology}

The syllabus phrase for this unit is: \textit{``Sociology -- The Discipline, Part A: Modernity and Social Changes in Europe and Emergence of Sociology.''} The instructor read this literally: the emergence of sociology is tied to social changes in Europe, and the throughline across all of them is modernity .

\subsection{The Middle Ages: The God-Centric World}
Before the Renaissance, in medieval Europe, the priestly class (the Church) dictated every aspect of a person's life---what to eat, whether one could travel, and even the shape of the universe .

\begin{KeyConcept}
\begin{itemize}
    \item God was at the centre of the universe; the world view was entirely religious .
    \item Individual ability or creativity was not acknowledged---painters could only paint godly/heavenly figures, not ordinary men, women or children .
    \item The Church held the view that the earth is at the centre of the universe (\textbf{geocentrism}) and that the human body is a ``divine product'' that must not be dissected or studied---this blocked the development of anatomy/medical science .
    \item Anyone who challenged the Church was declared a \textbf{heretic}---someone challenging its authority---and could be charged with \textbf{blasphemy} (insulting God or anything religious) .
    \item The aspiration of mankind in this period was to go to heaven, not any worldly achievement; ``life on earth is inconsequential'' was the Church's teaching, and one's afterlife depended on obedience to the clergy .
\end{itemize}
\end{KeyConcept}

\begin{QuickRevision}
\begin{itemize}
    \item \textbf{Middle Ages:} God-centric world view; Church controlled belief and knowledge; heresy/blasphemy suppressed challenge; individual creativity and anatomy/medicine were restricted .
\end{itemize}
\end{QuickRevision}

\subsection{The Renaissance: Rebirth of Man}
Renaissance (roughly 1400--1600) literally means rebirth/reawakening---rebirth of what man is capable of doing as an individual .

\begin{KeyConcept}
\begin{itemize}
    \item \textbf{Trigger:} Europeans came into contact with the Arab world through the Crusades (wars between Muslims and Christians) and trade, and through this contact rediscovered preserved Greek and Roman literature (Aristotle, Plato, the concepts of logic and reason) . Realising how advanced the Greeks and Romans had been, Europeans felt they had been kept underdeveloped by the priestly class .
    \item The Renaissance is characterised as the restoration of Greek and Roman literature and marks the departure from a God-centric to a man-centric world---mankind, not God, became the centre of the universe .
    \item \textbf{Philosophy of the Renaissance: Humanism}---centering on the human/mankind rather than God . Humanism encompasses (includes) both individualism and secularism .
\end{itemize}
\end{KeyConcept}

\begin{KeyConcept}
\begin{itemize}
    \item \textbf{Individualism:} Man's capacity to think and reason for himself, not simply accept what the priestly class dictates; the idea that individual human contribution and creativity can be celebrated (an idea previously reserved for God/the divine) .
    \item \textbf{Secularism (Renaissance-era sense):} Not against religion, but thinking beyond religion . Reason and faith co-existed rather than opposing each other in this period; the concept of God was not questioned outright---only supplemented by exploring human potential .
\end{itemize}
\end{KeyConcept}

\paragraph{Individual Contributions of the Renaissance}

\begin{ExampleBox}
\begin{itemize}
    \item \textbf{Nicolaus Copernicus (theory in 1490):} Proposed \textbf{heliocentrism} (from ``helio,'' root word for sun)---the sun, not the earth, is at the centre---directly countering the Church's geocentrism .
    \item \textbf{Johannes Kepler:} First to describe the laws of planetary motion, showing that planets move in an elliptical (not circular) orbit---contrary to the church-sanctioned (``Vatican law'') circular representation .
    \item \textbf{Galileo Galilei:} Verified Copernicus's claim through his own observation (built his own telescope) and openly opposed the Church; Copernicus, Kepler and Galileo were together declared heretics, charged with blasphemy, by the Church .
    \item \textbf{Leonardo da Vinci:} Called ``the Renaissance man''; studied human anatomy in defiance of Church restrictions; famous for the \textit{Mona Lisa} . His most famous painting was contrasted with pre-Renaissance religious art: earlier paintings of \textit{Madonna and Child} gave little attention to facial expression (glory was reserved for the divine), while Renaissance art (e.g., \textit{Mona Lisa}) used perspective to depict individual human emotion and expression---a marker of the shift from god-centric to man-centric art .
    \item \textbf{William Shakespeare:} Renowned for man-centred literary work (e.g., \textit{Hamlet}) about everyday human drama rather than purely religious/Biblical teaching, which had been the content of books in the Middle Ages .
    \item \textbf{Philippo Brunelleschi (named tentatively by the instructor):} Credited with the modern design of the printing press, significant because it allowed knowledge previously restricted to elite languages (Latin in Europe, Sanskrit in India, as opposed to the masses' Pali) to reach ordinary people .
    \item \textbf{Christopher Columbus} and the concept of \textbf{cartography} (map-making) and the ``age of exploration'' were cited as products of the same era---curiosity as a new trait of ``man'' .
\end{itemize}
\end{ExampleBox}

\begin{CriticalPoint}
\begin{itemize}
    \item Renaissance marked the departure from a God-centric world to a man-centric world . Its philosophy, humanism, revolves around regard and acknowledgement of the existence of the individual, and encompasses both individualism and secularism .
\end{itemize}
\end{CriticalPoint}

\begin{QuickRevision}
\begin{itemize}
    \item \textbf{Renaissance (1400--1600):} Humanism (individualism + secularism); departure from God-centric to man-centric world; Copernicus (heliocentrism), Kepler (elliptical orbits), Galileo, da Vinci, Shakespeare, the printing press .
\end{itemize}
\end{QuickRevision}

\subsection{The Scientific Revolution (Scientific Renaissance)}

\begin{KeyConcept}
\begin{itemize}
    \item The scientific revolution is described as another name for ``scientific renaissance''---not a separate period from the Renaissance but the strand of it in which man used reason, his own observation, and experimentation to understand the working of the (natural) universe .
    \item It is exemplified by the same figures already discussed---Copernicus, Kepler, Galileo---whose work was based on observation/experimentation rather than church teaching .
\end{itemize}
\end{KeyConcept}

\begin{GovtPolicy}
\begin{itemize}
    \item This period paved the way for the age of reason---the \textbf{Enlightenment era}, dated roughly 1600--1800 (with the Renaissance placed at roughly 1400--1600, and the scientific renaissance in between) .
\end{itemize}
\end{GovtPolicy}

\begin{QuickRevision}
\begin{itemize}
    \item \textbf{Scientific Revolution = ``scientific renaissance'':} Reason/observation/experimentation applied to the natural universe (Copernicus, Kepler, Galileo); paved the way for the Enlightenment .
\end{itemize}
\end{QuickRevision}

\subsection{The Enlightenment: The Age of Reason}
Enlightenment (``age of reason'') again involves man using reason, observation and experimentation---but now applied not to the natural universe but to improve the working of society, meaning specifically the reform of governance .

\begin{KeyConcept}
\textbf{Distinction drawn:} In the scientific revolution, non-human/natural phenomena were the object of inquiry; in the Enlightenment, man himself became the object of inquiry to other men---scholars studied ``what is the nature of man'' in order to work out how governance could be improved .

The instructor described Enlightenment thinkers as elites who met and debated in \textbf{salons} (informal gatherings/cafes) in France---compared to modern-day coffee-shop discussions among the educated .
\end{KeyConcept}

\begin{QuickRevision}
\begin{itemize}
    \item \textbf{Enlightenment (1600--1800):} Reason applied to reform of governance/society (not the natural universe); for the first time, man himself becomes the object of study .
\end{itemize}
\end{QuickRevision}

\subsection{Civil Society}

\begin{KeyConcept}
\textbf{Civil society}---defined as anything that falls between family and state: media (print, digital, social), research institutions, NGOs, pressure groups, university professors/academicians, scientists, resident welfare associations (RWAs), and ordinary individuals critical of government on social media .
\begin{itemize}
    \item \textbf{Function of civil society:} To keep a check and balance on the power of the state .
    \item The concept of civil society is traced to the same French salons where Enlightenment elites first began speaking against/critiquing governance---the instructor draws a direct line from those salons to present-day newspaper editorials (\textit{The Hindu}, \textit{The Indian Express}) and social media commentary critiquing government policy .
\end{itemize}
\end{KeyConcept}

\begin{QuickRevision}
\begin{itemize}
    \item \textbf{Civil society:} Anything between family and state (media, NGOs, research bodies, academics); function is to check and balance state power; traced to the French salons .
\end{itemize}
\end{QuickRevision}

\subsection{Enlightenment Scholars and the Critique of Monarchy}
All Enlightenment scholars discussed were concerned with one shared question: how can governance (and later, economy and gender relations) be improved through the application of reason . Each is discussed below as engaging with the French monarchical system of the time .

\paragraph{Thomas Hobbes}
\begin{KeyConcept}
\textbf{Context:} Wrote during the English civil wars; his book is \textit{Leviathan} .
\begin{itemize}
    \item \textbf{Hobbes's view of human nature:} Men are inherently lazy, stupid, greedy and selfish .
    \item Without government, Hobbes argued, life would be \textit{``solitary, poor, nasty, brutish and short''} (quoted by the instructor as Hobbes's own formulation) .
    \item \textbf{His prescription:} A strong monarch is needed to keep these impulses in check and maintain social order .
    \item \textbf{Important nuance:} Although Hobbes supported monarchy, he did \textit{not} support the divine right of rule (i.e., the idea that a monarch's authority comes from God/a lineage sanctioned by God)---which is why he is still counted among Enlightenment (not purely traditionalist) thinkers .
\end{itemize}
\end{KeyConcept}

\paragraph{John Locke}
\begin{KeyConcept}
\begin{itemize}
    \item \textbf{Locke's view of human nature:} Men are not stupid or selfish but moral beings, possessing natural rights to \textbf{life, liberty and property} .
    \item Because no one, in Locke's time, held private property (land ownership was in the name of the crown/king, who distributed it only to nobles, not the masses), Locke's argument was in effect that people were being deprived of a right they should naturally have .
    \item Locke held that the duty of the monarch is only to protect these natural rights (life and property)---this exchange (people accept government in return for protection of natural rights) is the origin of the idea of a \textbf{social contract} between men and the state .
\end{itemize}
\end{KeyConcept}

\begin{GovtPolicy}
\textbf{Social contract explained by analogy:} Paying taxes (direct or indirect) to government in exchange for government action (infrastructure, public goals) is a modern-day form of the same bargain---citizens give up a portion of freedom in exchange for social order and protection; this is why fundamental rights are conditional, not absolute .

\textbf{Hobbes vs Locke on revolt:} Hobbes held the monarch has full right to crush any revolt; Locke held that if the monarch fails to protect people's rights, the people have a right to revolt---the monarch cannot suppress a revolt born of failing to honour the social contract .
\end{GovtPolicy}

\paragraph{Jean-Jacques Rousseau}
\begin{KeyConcept}
\textbf{Concept: General will}---the will of the majority; directly linked by the instructor to the modern concept of democracy .
\begin{itemize}
    \item \textbf{Rousseau's view of human nature:} Man is essentially wise and good, ``born free,'' but it is society that puts him into chains---illustrated by the idea that a newborn has no concept of nationality, caste, class, sexuality, religion, ethnicity or race, but acquires and becomes bound by all of these through socialisation .
    \item Rousseau also wrote about the rights of women and women's education, but the instructor flagged this as a reflection of a patriarchal mindset: Rousseau's proposed education for women was oriented toward being a good housewife, not toward the fuller rights he proposed for men .
\end{itemize}
\end{KeyConcept}

\paragraph{Voltaire}
\begin{KeyConcept}
\begin{itemize}
    \item \textbf{Position:} Called for religion to be ``crushed''---meaning that there should be no exploitation of man in the name of religion/God, since man cannot empirically verify God and has, in the instructor's words, ``been fooled for thousands of years'' by rule justified through religion .
    \item Also a strong supporter of freedom of speech, with the paraphrased line attributed to him: he may disagree with what someone says, but would defend, to the point of death, their right to say it .
\end{itemize}
\end{KeyConcept}

\begin{ExampleBox}
The instructor connected Voltaire's critique of religiously justified hierarchy to the Hindu Varna system (Brahmins at the top, followed by Kshatriyas, Vaishyas, Shudras) and named the Buddha as the historical Indian figure who first challenged this order---noting that the Hindi word for enlightenment, \textit{Nirvana}, derives from this association .
\end{ExampleBox}

\paragraph{Montesquieu (Baron de Montesquieu)}
\begin{KeyConcept}
\begin{itemize}
    \item \textbf{Position:} Governance improves when power is not concentrated in one person---under absolute monarchy, one king held legislative power (made laws), executive power (implemented laws, appointed anyone without merit-based selection), and judicial power (acted as judge) all at once, which the instructor calls despotism/autocracy .
    \item \textbf{Proposed solution: Separation of powers}---three distinct branches of government: legislature (formulates laws through debate/deliberation, e.g., Parliament), executive (implements laws once passed as acts), and judiciary .
    \item The instructor added a practical note for GS/polity: in practice an IAS officer may act across all three functions at different times (e.g., a District Magistrate exercising judicial powers), so the separation is not always a watertight compartment in practice .
\end{itemize}
\end{KeyConcept}

\paragraph{Adam Smith}
\begin{KeyConcept}
Adam Smith is described as part of the Scottish Enlightenment, itself part of the broader European Enlightenment . His book: \textit{Wealth of Nations} .
\begin{itemize}
    \item \textbf{Laissez-faire}---``let it be free''---a free-market economy not intervened in by government; linked to the modern term \textit{neoliberalism}, meaning expansion of free markets with minimum government role . The instructor illustrated with import duty: a fruit that should cost \rupee 50 costs \rupee 200 after import duty, pushing consumers toward ``indigenous'' goods---he notes Smith would not approve, since protectionism (tariffs, farm subsidies) works against free-market principles .
    \item \textbf{Invisible hand}---self-interested individuals (e.g., entrepreneurs seeking profit) unintentionally benefit society as a whole . Illustrated with Coca-Cola vs Pepsi: competing to capture market share by lowering prices benefits consumers (choice, lower prices) and creates employment---benefits arise without either firm intending to serve society directly, and only where government does not intervene .
    \item \textbf{Post-Fordism} (mentioned in passing, to be developed later in the syllabus under work/economy): an era where products are customised to individual consumer preference (e.g., a custom-assembled computer, a personalised printed t-shirt or mug) .
\end{itemize}
\end{KeyConcept}

\begin{CriticalPoint}
\textbf{Herbert Spencer}, described as a ``social Darwinist,'' applied laissez-faire thinking to society itself: no protection or subsidy should be given to the poor; only the fittest should survive . The instructor stressed that the phrase ``survival of the fittest'' was coined by Spencer, not by Charles Darwin, though Spencer was influenced by both Darwin and Adam Smith . The term ``social Darwinist project'' is used in media commentary (e.g., \textit{The Hindu}, \textit{The Indian Express}) to describe situations where government does not protect people (e.g., during crises such as lockdowns), leaving survival to chance .
\end{CriticalPoint}

\paragraph{Mary Wollstonecraft}
\begin{KeyConcept}
\textbf{Book:} \textit{A Vindication of the Rights of Woman} (also referenced by the instructor as \textit{A Serious Proposal to the Ladies}) .
\begin{itemize}
    \item \textbf{Position:} Society consists equally of women, and if men can read, write and study, so can women . She is described by the instructor as the mother/founder of feminism---the first Enlightenment scholar to speak assertively, exclusively, about women's rights and women's education (as distinct from Rousseau's more limited, housewife-oriented proposal) .
\end{itemize}
\end{KeyConcept}

\begin{PYQLink}
A warm-up matching exercise was given, asking students to match each of the following ideas to its scholar :
\begin{enumerate}[label=(\roman*)]
    \item \textit{Leviathan} / social contract / ``life would be solitary, poor, nasty, brutish and short'' without government $\longrightarrow$ \textbf{Hobbes} .
    \item Improving the status of women, ``if all men are born free, how is it that all women are born slaves?'' $\longrightarrow$ \textbf{Wollstonecraft} .
    \item Civilization corrupts, man is born free but everywhere in chains, government improved through the general will $\longrightarrow$ \textbf{Rousseau} .
    \item A French writer who studied political liberty, was impressed by the British system, and held that ``power should be a check to power'' $\longrightarrow$ \textbf{Montesquieu} .
\end{enumerate}
\end{PYQLink}

\begin{QuickRevision}
\begin{itemize}
    \item \textbf{Hobbes:} Monarchy, men are stupid/selfish, rejects divine right .
    \item \textbf{Locke:} Natural rights, social contract, right to revolt .
    \item \textbf{Rousseau:} General will/democracy, man born free .
    \item \textbf{Voltaire:} Anti-clericalism, freedom of speech .
    \item \textbf{Montesquieu:} Separation of powers .
    \item \textbf{Adam Smith:} Laissez-faire, invisible hand, post-Fordism (Spencer/social Darwinism nearby) .
    \item \textbf{Wollstonecraft:} Women's rights, founder of feminism .
\end{itemize}
\end{QuickRevision}

\subsection{Feudal Land Ownership: Europe and India Compared}

\begin{GovtPolicy}
\begin{itemize}
    \item In feudal Europe, land was owned by the crown (the king), not by individuals; a noble who received land from the king had no hereditary right to pass it to his son, so there was little incentive to invest in or improve the land---land therefore stayed unproductive .
    \item \textbf{Comparison made to India:} During the Mughal era there was likewise no concept of private land ownership---land ownership was, similarly, in the name of the crown .
    \item \textbf{Scholar cited: François Bernier}, a physician-traveller (not a historian or sociologist) who visited India during the reign of the Mughal emperor Aurangzeb and drew the comparison between feudal Europe and feudal India .
    \item Private land ownership arrived in India only under British rule, which introduced the \textbf{Zamindari system}; the earlier \textbf{Jagirdari system} had existed under the Sultanate (before the Mughals) . \textbf{Difference:} Zamindari was hereditary (a zamindar could pass land to his son and collected revenue for the Empire); Jagirdari had no such hereditary right .
    \item Because Zamindari lacked strong incentive for investment (among other reasons), productivity stayed low, contributing---alongside the same dynamic in France---to conditions of famine .
\end{itemize}
\end{GovtPolicy}

\begin{QuickRevision}
\begin{itemize}
    \item Feudal Europe and Mughal India both had crown-owned (not private) land (François Bernier's comparison); India got private ownership only under British rule (hereditary Zamindari; earlier non-hereditary Jagirdari under the Sultanate) .
\end{itemize}
\end{QuickRevision}

\subsection{The French Revolution: Background and the Estates System}
The instructor framed the ideas of Enlightenment scholars as reforms proposed against the actual French monarchical system they lived under, and set out the structure of pre-revolutionary French society, which was divided into three ``estates'' .

\vspace{6pt}
\noindent
\begin{tabularx}{\textwidth}{l p{4.2cm} c c Y}
\toprule
\rowcolor{AccentDark}
\thead{Estate} & \thead{Composition} & \thead{\% Pop.} & \thead{\% Land} & \thead{Taxes Paid} \\
\midrule
\textbf{First Estate} & Catholic Church officials and clergy & 1\% & 10\% & None  \\
\textbf{Second Estate} & Nobility and aristocracy (earls, dukes, knights, barons) & 2\% & 20\% & None  \\
\textbf{Third Estate} & Bourgeoisie (educated city lawyers, doctors, businessmen) and peasants/farmers & 97\% & 70\% & All taxes, plus tithes (unevenly distributed; many peasants landless)  \\
\bottomrule
\end{tabularx}
\vspace{6pt}

\begin{DataFact}
\begin{itemize}
    \item Together, the first and second estates---3\% of the population---owned 30\% of the land and paid no tax; the third estate---97\% of the population---held the remaining 70\% of the land (very unevenly, with many peasants entirely landless) and bore the full tax burden .
    \item \textbf{Tithe}---a mandatory contribution of one-tenth ($1/10$) of one's produce or earnings to the Church, paid by ordinary people (e.g., a farmer's harvest, a businessman's earnings) .
    \item A cartoon of the period was described: the third estate depicted as a horse carrying the burden of the first and second estates seated upon it .
\end{itemize}
\end{DataFact}

\begin{CriticalPoint}
\begin{itemize}
    \item The third estate was described in the lecture as ``everything'' in terms of population and productive/tax contribution, yet had ``nothing'' in terms of political voice, and ``desired to be something''---the instructor named \textbf{Abbé Sieyès} as the figure who inspired the third estate to seek political power .
\end{itemize}
\end{CriticalPoint}

\paragraph{Economic and Political Causes}

\begin{DataFact}
\textbf{Economic causes cited:}
\begin{enumerate}[label=(\roman*)]
    \item \textbf{Extravagant royal spending}---termed ``Madame Deficit,'' referring to the fiscal profligacy attributed to Queen Marie Antoinette's extravagant expenses on exotic fruits, clothes and jewellery (the instructor noted this framing itself reflects a patriarchal tendency to single out the Queen for blame) .
    \item \textbf{The Seven Years' War} between France and Britain for global supremacy, which drained the treasury .
    \item \textbf{Famine}---soldiers went hungry even during the war .
\end{enumerate}
\textbf{Political context:} King Louis XVI is described as an absolute monarch but a weak ruler who attempted reforms that mostly failed due to resistance from the nobility, and who incurred huge debt partly linked to the palace of Versailles .
\end{DataFact}

\begin{ExampleBox}
A period proverb quoted by the instructor: \textit{``The king in his castle, the poor at his gate; God made them highly and lowly, and ordered their estates''}---used to show how religious framing (``God has ordered these estates'') discouraged any challenge to the existing order, in the way the instructor said caste hierarchy was maintained in India ``until Ambedkar voiced out their opinion'' .
\end{ExampleBox}

\begin{QuickRevision}
\begin{itemize}
    \item \textbf{French estates:} 1st (clergy, 1\% pop., 10\% land, no tax), 2nd (nobility, 2\% pop., 20\% land, no tax), 3rd (bourgeoisie + peasants, 97\% pop., 70\% land, all taxes + tithes) .
    \item \textbf{Causes:} Enlightenment ideas + Madame Deficit + Seven Years' War + famine .
\end{itemize}
\end{QuickRevision}

\subsection{The French Revolution: Estates General to Republic}

\begin{GovtPolicy}
\begin{itemize}
    \item \textbf{Estates General}---the rarely-convened body (equivalent to a parliamentary session) in which all three estates met . Each estate had only one vote regardless of population, so the first and second estates (allied with each other) always outvoted the third estate 2-to-1, despite representing only 3\% of the population .
    \item \textbf{Tennis Court Oath}---members of the third estate broke away and formed their own \textbf{National Assembly}, a direct challenge to the monarch and the existing hierarchical/ruling structure, asserting the right to make their own rules .
    \item The National Assembly issued the \textbf{Declaration of the Rights of Man and of the Citizen}, proposing that everyone should have the right to life, property and vote---though the instructor noted that in practice women and slaves were still excluded from voting and not considered active citizens .
\end{itemize}
\end{GovtPolicy}

\begin{GovtPolicy}
\begin{itemize}
    \item \textbf{First proposal accepted by the National Assembly: ``Law of clergy,''} i.e., a secular state---the monarch/King Louis's family would no longer rule in the name of God; religion and rule were to be separated .
    \item \textbf{Outcome of the Revolution:} Transition from absolute monarchy to a republic . The instructor noted this transition was not immediate or smooth---it took a long time for a stable republic to emerge, and the French Revolution is understood by the instructor to have recurred more than once (1789/1791 and again in the 1800s), though it is usually studied as a single 1790s event .
\end{itemize}
\end{GovtPolicy}

\paragraph{What the French Revolution Contributed to Modernity}

\begin{KeyConcept}
\begin{enumerate}
    \item \textbf{Democracy}---ideas on government (rights to life, liberty, property; Rousseau's general will; Montesquieu's separation of power) directly inspired calls for democratic government .
    \item \textbf{Communism (the seed of it)}---the National Assembly proposed that all land and institutions (banks etc.) confiscated from the first/second estates be nationalised and equally redistributed rather than concentrated as before; the instructor called this the advent of communism/an egalitarian society, while noting that the boundary between socialism and communism is highly blurred and would be discussed later .
    \item \textbf{Nationalism}---a group of people bound by allegiance to a common thought process/common codes can be considered a ``nation,'' which need not correspond to a ``state'' (territory + political administration) . The third estate coming to see itself as a unified nation, having been ``united'' in shared exploitation, is presented as the origin of modern nationalism; the instructor drew a parallel to Indian nationalism as, similarly, a by-product of the shared experience of colonialism .
    \item \textbf{Citizenship}---the idea that all persons, once made citizens, hold equal rights irrespective of any prior estate, caste, lineage, gender, religion or race (the instructor cited Article 14 of the Indian Constitution as a present-day parallel) .
\end{enumerate}
\end{KeyConcept}

\begin{QuickRevision}
\begin{itemize}
    \item Estates General (one vote per estate) $\rightarrow$ Tennis Court Oath $\rightarrow$ National Assembly $\rightarrow$ Declaration of the Rights of Man $\rightarrow$ secular state/republic .
    \item French Revolution gave modernity: democracy, the seed of communism, nationalism, and citizenship .
\end{itemize}
\end{QuickRevision}

\subsection{The Idea of Modernity: A Consolidated Model}
The instructor drew a single flowchart :
$$\text{Middle Ages} \longrightarrow \text{Renaissance} \longrightarrow \text{Scientific Renaissance} \longrightarrow \text{Enlightenment} \longrightarrow \text{French Revolution} \longrightarrow \text{Industrial Revolution} \longrightarrow \text{Modernity}$$
Social changes in Europe are the cause; modernity is their consequence .

\begin{KeyConcept}
\begin{itemize}
    \item \textbf{Modernity is defined as the decline (not death) of traditional beliefs, ideas and institutions}---those institutions weaken but do not disappear entirely .
    \item \textbf{René Descartes} is introduced here as a philosopher straddling the scientific revolution/Enlightenment: his statement \textit{``the only thing that cannot be doubted is doubt itself''} is read as establishing that man's capacity to think/doubt proves his own existence---leading to his famous line, \textit{``I think, therefore I am''} (\textit{Cogito, ergo sum}), directly challenging the Catholic Church's view that only the soul, not the individual man, exists .
\end{itemize}
\end{KeyConcept}

\begin{KeyConcept}
The consolidated ``blueprint'' of the elements of modernity, as drawn by the instructor from all events studied, was given as (not in strict chronological order) :
\begin{enumerate}
    \item \textbf{Humanism} (philosophy of the Renaissance), encompassing individualism and secularism 
    \item \textbf{Rationality}---man is a rational, thinking being, not merely a soul awaiting the afterlife 
    \item \textbf{Empiricism}---belief only in what can be sensed (seen, heard, touched); the basis of science 
    \item \textbf{Science}---rationality and empiricism together 
    \item \textbf{Reason} (from the Enlightenment) 
    \item \textbf{Democracy} 
    \item \textbf{Freedom of speech / expression} 
    \item \textbf{Gender equality / feminism} (from Wollstonecraft) 
    \item \textbf{Nationalism} (from the French Revolution) 
    \item \textbf{Citizenship} 
    \item \textbf{Republic / the state} 
    \item \textbf{Nuclear family} (an unintended consequence of the Industrial Revolution, discussed below) 
    \item \textbf{Industrial economy / urbanisation / social class} as the primary modern identity (replacing estate, caste, or lineage) 
\end{enumerate}
\end{KeyConcept}

\begin{CriticalPoint}
\begin{itemize}
    \item The instructor stressed there is no fixed, universally agreed number of elements of modernity (``it all depends from book to book'')---what matters is the overall structure and that all of these elements recur across later units of the syllabus (politics and society, family and kinship, religion, work and economy) .
\end{itemize}
\end{CriticalPoint}

\begin{QuickRevision}
\begin{itemize}
    \item \textbf{Modernity = decline (not death) of traditional institutions} .
    \item Descartes' ``I think, therefore I am'' as the philosophical marker of individual existence .
    \item Modernity's 13-point blueprint spans humanism, rationality, empiricism, science, reason, democracy, free speech, gender equality, nationalism, citizenship, republic, nuclear family, and industrial economy/urbanisation/class .
\end{itemize}
\end{QuickRevision}

\subsection{The Industrial Revolution}

\begin{KeyConcept}
\textbf{Industrial Revolution:} A period, beginning around 1750 in England, marked by the shift from man-made (hand/craft) production to machine-made, mass production . Unlike the other events discussed, the instructor stresses it has no clear end date---it is ongoing (``Industrial Revolution 4.0'' today) .

\textbf{Cottage industry}---small-scale, home-based production in which family members together manufacture goods, illustrated with beedi-making and Lijjat papad, both largely made by women at home .
\begin{itemize}
    \item The instructor tied the possibility of industrial invention (steam engine, electric bulb, the Wright brothers' first flight, Thomas Alva Edison's inventions) back to the earlier chain of events: none of this would have been possible without the scientific revolution, which in turn would not have been possible without the Renaissance .
\end{itemize}
\end{KeyConcept}

\subsection{Pre-Industrial vs Post-Industrial Society}
The instructor set out a point-by-point comparison, noting it helps understand not just the technological changes brought by the Industrial Revolution but the social changes that followed from them .

\vspace{6pt}
\noindent
\begin{tabularx}{\textwidth}{p{3.2cm} Y Y}
\toprule
\rowcolor{AccentDark}
\thead{Feature} & \thead{Pre-Industrial Revolution} & \thead{Post-Industrial Revolution} \\
\midrule
\textbf{Economic system} & Agrarian economy / feudalism (land as primary source of livelihood) & Industrial economy / capitalism (mass production through factories/machines)  \\
\textbf{Scale of industry} & Cottage industry, small scale (home-based, e.g., beedi, Lijjat papad) & Heavy-scale industry (steel, big dams, assembly lines)  \\
\textbf{Division of labour} & Simple division of labour (tasks divided informally within the family) & Complex division of labour, based on specialisation/degrees/skills, requiring integration of hundreds of specialised workers  \\
\textbf{Source of power} & Water mills / wind mills & Coal and steam (steam engine allowed factories to run all day and night)  \\
\textbf{Economic motive} & Subsistence economy (production for one's own livelihood, not profit) & Profit-centric economy (mass production/consumption for commercial profit)  \\
\textbf{Home and workplace} & Unified---home was the workplace & Separated---workplace (factory) distinct from home; this separation was also sexual in nature (home increasingly managed by women, factory work increasingly sought by men)  \\
\textbf{Family structure} & Joint family system & Nuclear family system (linked to rural-to-urban migration/urbanisation for job opportunities)  \\
\bottomrule
\end{tabularx}
\vspace{6pt}

\begin{CriticalPoint}
\begin{itemize}
    \item The instructor connected the specialisation point back to \textbf{alienation}: a man who might once have made an entire car by himself now, under complex division of labour, only makes a single nut for the wheel throughout his working life---specialisation, in this reading, has ``killed'' man's ability to explore his fuller potential .
\end{itemize}
\end{CriticalPoint}

\begin{QuickRevision}
\begin{itemize}
    \item \textbf{Pre $\rightarrow$ Post Industrial shifts:} Agrarian/feudal $\rightarrow$ industrial/capitalist; cottage $\rightarrow$ heavy industry; simple $\rightarrow$ complex division of labour; water/wind $\rightarrow$ coal/steam; subsistence $\rightarrow$ profit economy; unified $\rightarrow$ separated home/work; joint $\rightarrow$ nuclear family .
\end{itemize}
\end{QuickRevision}

\subsection{Social Evils of the Industrial Revolution and Marx's Response}
Alongside its benefits (urban expansion, job opportunities, railways, aeroplanes), the instructor listed several ``social and economic costs'' of the Industrial Revolution, urging students not to treat these as purely individual failings but as structural consequences .

\begin{CriticalPoint}
\begin{itemize}
    \item \textbf{Slum development}---poverty around industrial corridors framed as a structural issue, not an individual family's failing; giving money or books to a child begging is, in this reading, treating a structural problem as an individual one .
    \item \textbf{Class inequality}---widened because only some could access new opportunities (education, mobility), while others remained constrained by family circumstance, gender, religion or caste .
    \item \textbf{Overcrowding of urban centres} (illustrated with the example of Gurgaon) .
    \item \textbf{Industrial pollution} .
    \item \textbf{Exploitation of workers}---with far more job-seekers than available jobs, a ``reserve army'' of unemployed people became available as cheap labour, and those who were employed lived in fear of being replaced .
    \item \textbf{Degradation of work}---a traditional craftsperson (e.g., a potter) who once had full control over the shape, quality and quantity of what they made lost that control once absorbed into industrial production, where supervisors dictate output and a single item passes through many specialised hands; the traditional guild system (specialised craft streets/areas) declined as machine-made goods, cheaper and branded, out-competed hand-made goods .
\end{itemize}
\end{CriticalPoint}

\begin{GovtPolicy}
Karl Marx wrote the \textit{Communist Manifesto} in response to these unintended, exploitative consequences of the Industrial Revolution---workers who once had full control over their labour were now treated as cheap labour, working long hours with no concept of minimum wage or fixed working hours (unlike today's standards) .

The \textit{Manifesto}'s closing call, paraphrased by the instructor: workers of the world should unite, since they have nothing to lose but their (exploitative) chains, and should overthrow the industrialist/bourgeois class .

Marx is also described as calling the French Revolution itself a ``bourgeois revolution''---in his reading, it was the bourgeoisie (not the peasantry) within the third estate who ultimately gained the most from it .
\end{GovtPolicy}

\begin{QuickRevision}
\begin{itemize}
    \item \textbf{Social evils of industrialisation:} Slums, class inequality, overcrowding, pollution, worker exploitation, degradation of craft work---Marx responded with the \textit{Communist Manifesto}, calling workers to unite and overthrow the bourgeoisie .
\end{itemize}
\end{QuickRevision}

\subsection{Conservative Reaction to Modernity}
The instructor introduced this section by asking whether changes such as the shift from joint to nuclear family, or the arrival of secularism, are ``good'' or ``bad''---stressing that in sociology there is no absolute answer: what is celebrated by one perspective is opposed by another, depending on one's social location and ideology .

\begin{KeyConcept}
\textbf{Counter-Enlightenment}---a period/school of thought that emerged in response to (against) the ideas of the Enlightenment, the French Revolution and the Industrial Revolution .
\begin{itemize}
    \item Conservatives believed modernity undermines (downgrades the importance of) traditional beliefs and institutions such as patriarchy, religion/the Church, monogamous marriage, and monarchy, and wanted to restore that earlier order .
\end{itemize}
\end{KeyConcept}

\begin{CriticalPoint}
\textbf{Why conservatives wanted to restore the old order:} The period after the French Revolution was not orderly---it was marked by chaos, turmoil, and what the instructor (borrowing from Durkheim) called \textbf{normlessness} . The Reign of Terror (c.~1794 onward) saw mass executions of those who had opposed the Revolution, including King Louis XVI and Marie Antoinette themselves .
\begin{itemize}
    \item \textbf{Conservatives argued:} If this chaos and bloodshed is what modernity/democracy brings, is it worth celebrating?  They wanted a return to a time of ``proper order,'' even though (in the instructor's own assessment) a society, like an adult, cannot simply be returned to its earlier/``infant'' stage .
    \item Conservatives were also unhappy with expanded women's rights (e.g., the right to divorce, gained after the French Revolution) and framed rising divorce rates as a threat to family as an institution; the instructor illustrated the underlying tension with feminist theorising of a parallel ``sexual contract'' between men and women (paralleling Hobbes/Locke/Rousseau's social contract between men and the state) once work in the home came to be seen as assigned by sex .
\end{itemize}
\end{CriticalPoint}

\paragraph{Key Conservative Scholars Named}

\begin{GovtPolicy}
\begin{itemize}
    \item \textbf{Louis de Bonald and Joseph de Maistre}---described as extreme conservatives, i.e., those who wanted to fully restore the pre-modern order (monarchy, patriarchy, Church authority) as it had existed---which the instructor compares to asking an adult to literally become a child again .
    \item \textbf{Auguste Comte}---named as conservative despite being strongly associated with science; the instructor stressed science and conservatism are not opposites (citing Einstein's line, paraphrased: science without religion is ``lame,'' religion without science is ``blind'') .
    \item \textbf{Emile Durkheim}---named explicitly in the syllabus as conservative; discussed in detail below .
    \item \textbf{Henri de Saint-Simon}---named alongside Comte as another conservative-leaning figure of this transitional period .
\end{itemize}
\end{GovtPolicy}

\begin{CriticalPoint}
\textbf{Durkheim's specific discontents (as presented):} He was depressed by several forms of social disorder :
\begin{enumerate}[label=(\roman*)]
    \item \textbf{Anti-Semitism} (hatred against Jews, which he observed being carried out amid the same era's wars, including the Franco-Prussian war, and which the instructor also links forward to Hitler's targeting of Jews in Germany) ;
    \item \textbf{Church-state discord}, i.e., secularism/the separation of church and state, which he was not happy with ;
    \item \textbf{Urbanization}, because it weakens family ties---illustrated with the example of a child moving abroad (e.g., into the Indian Foreign Service) and the family's mixed feelings of pride and quiet unhappiness at the resulting distance .
\end{enumerate}
\end{CriticalPoint}

\begin{QuickRevision}
\begin{itemize}
    \item \textbf{Counter-Enlightenment/conservatives (de Bonald, de Maistre = extreme; Comte, Durkheim, Saint-Simon):} Sought restoration of tradition via reform, alarmed by post-Revolution disorder (Reign of Terror); Durkheim specifically opposed anti-Semitism, church-state separation, and urbanisation's effect on family .
\end{itemize}
\end{QuickRevision}

\subsection{The Radical Reaction: Karl Marx}

\begin{KeyConcept}
Where conservatives sought a return to the pre-modern order through reform, Marx is presented as the radical reaction: he held that workers should overthrow the industrialist/bourgeois class outright, and that meaningful social change comes only through revolution, not gradual reform .
\begin{itemize}
    \item Marx read the French Revolution itself as a bourgeois revolution and, more broadly, urged workers to revolt against the bourgeoisie who, in his reading, exploited them without paying what they were truly owed .
\end{itemize}
\end{KeyConcept}

\begin{CriticalPoint}
\textbf{Conservatives vs radicals on why they differ:}
\begin{itemize}
    \item \textbf{Conservatives} believe revolution creates only more disorder (``how much more disorder do you want, after you've already killed the king, queen, nobility?'') and that social reform, not revolution, restores order .
    \item \textbf{Radicals (Marxists)} hold that reform is a way of preventing (delaying/defusing) revolutionary change rather than truly resolving exploitation---illustrated by an example of a company, police and government acting in a ``tight nexus'' (``cash nexus'') to pacify a labour strike with promised talks rather than addressing exploitation .
\end{itemize}
\end{CriticalPoint}

\begin{QuickRevision}
\begin{itemize}
    \item \textbf{Radicals (Marx):} Social change comes only through revolution, not reform; reform is seen as delaying real change (``cash nexus'' of state, police and capital) .
\end{itemize}
\end{QuickRevision}

\subsection{The Emergence of Sociology}
Having covered the elements of modernity and the two broad reactions to it (conservative and radical), the instructor closed the unit by explaining how sociology itself emerged from this context .

\begin{KeyConcept}
\begin{itemize}
    \item \textbf{Structural functionalists}---those who, in response to the problems of modernity, believed change should come via reform, not revolution . The instructor's analogy: society evolves like a natural human body; when the body/society is hurt, the response is to treat and heal (reform), not to hunt down and eliminate whoever caused the injury (revolution) .
    \item \textbf{Marxists}---the radicals, discussed above, who reject reform in favour of revolutionary change; the instructor's stated reason for their scepticism of reform: in the name of reform, revolutionary change is prevented/deferred rather than truly delivered .
\end{itemize}
\end{KeyConcept}

\begin{CriticalPoint}
\textbf{Sociology emerged, in the instructor's framing, for two purposes:}
\begin{enumerate}
    \item \textbf{To find the social laws that can restore order in society}---just as physics has laws like the law of gravity or planetary motion, sociologists (``scientists of society'') sought equivalent laws governing society ; and
    \item \textbf{To study the social changes and social evils associated with modernity} .
\end{enumerate}

\textbf{Positivism}---introduced briefly as the philosophy/method through which social laws governing society can be discovered, associated here specifically with \textbf{Auguste Comte} . The instructor described positivists as analogous to meteorologists who predict weather, or as scholars who might, in principle, be able to predict geopolitical outcomes (e.g., when a conflict will resolve) once the underlying ``laws'' governing a society are known .

Both structural functionalism and Marxism are described as positivist philosophies in this specific sense: both are engaged in trying to discover laws that govern the movement of society, even though they differ sharply on whether reform or revolution is the right response to those laws .
\end{CriticalPoint}

\begin{CriticalPoint}
\textbf{The instructor's closing reflection:} Having discovered ``man'' in the Renaissance and then set man to studying the natural universe (scientific revolution) and to reforming governance (Enlightenment), sociology represents a further, self-referential step---\textbf{man studying man}---to the point where the instructor remarks that the sociologist studying society risks becoming a kind of new authority (``new God'') in his or her own right, and that for committed Marxists specifically, Marx himself can function almost as a quasi-religious authority (``Marxism is a religion'' in that sense) .
\end{CriticalPoint}

\begin{QuickRevision}
\begin{itemize}
    \item Sociology emerged to find social laws (positivism, Comte) that could restore social order and to study modernity's social evils; structural functionalism (reform) and Marxism (revolution) are both, in this sense, positivist philosophies .
\end{itemize}
\end{QuickRevision}

\sectiondivider

\section{Full-Topic Revision Sheet}

\subsection{Topic 1 Revision --- The Sociological Perspective}

\begin{QuickRevision}
\begin{itemize}
    \item \textbf{Sociological perspective:} Seeing the system/force/structure behind individuals, not just the individuals themselves (solar system, gravity, benzene analogies)[cite: 1, 2].
    \item \textbf{Conspicuous consumption:} Buying to display status, not for use-value (Starbucks example).
    \item \textbf{Poverty/child labour read structurally:} Inequality, unemployment, malnutrition, cheap labour---not individual failing.
    \item \textbf{Division of labour (Adam Smith, \textit{Wealth of Nations})} vs \textbf{sexual division of labour} (task assignment by gender, not skill).
    \item \textbf{Media reproduces patriarchy;} \textbf{alienation} and \textbf{emotional labour} (Marx) describe the loss of control/authenticity in modern work.
    \item \textbf{Gender and occupation:} Segregation by sex in imagined roles (army officer); gated communities show voluntary class segregation.
    \item \textbf{Hidden curriculum:} School teaches authority/discipline/conformity beyond the official syllabus; surveillance sustains conformity and deters deviance.
    \item \textbf{Caste stratification (India)} compared to the French estates system as a bridge to the historical unit; \textbf{heteronormativity} and \textbf{deviance} introduced via the family picture.
\end{itemize}
\end{QuickRevision}

\subsection{Topic 2 Revision --- Modernity, Social Changes in Europe \& Emergence of Sociology}

\begin{QuickRevision}
\begin{itemize}
    \item \textbf{Middle Ages:} God-centric, Church-controlled world view; heresy/blasphemy charges suppressed challenge.
    \item \textbf{Renaissance (1400--1600):} Humanism = individualism + secularism; Copernicus (heliocentrism), Kepler (elliptical orbits), Galileo, da Vinci, Shakespeare.
    \item \textbf{Scientific Revolution = ``scientific renaissance'':} Reason/observation/experimentation applied to the natural universe; paved the way for the Enlightenment.
    \item \textbf{Enlightenment (1600--1800):} Reason applied to reform of governance/society; man becomes object of study; civil society originates in salons.
    \item \textbf{Thinkers compared:} Hobbes (monarchy, men are stupid/selfish, no divine right) vs Locke (natural rights, social contract, right to revolt) vs Rousseau (general will/democracy, man born free) vs Voltaire (anti-clericalism, free speech) vs Montesquieu (separation of powers) vs Adam Smith (laissez-faire, invisible hand, post-Fordism; Spencer/social Darwinism nearby) vs Wollstonecraft (women's rights, founder of feminism).
    \item \textbf{Feudal land:} Crown-owned in both feudal Europe and Mughal India (François Bernier); private ownership in India arrived only under British rule (Zamindari, hereditary; earlier Jagirdari, non-hereditary).
    \item \textbf{French Estates:} 1st (clergy, 1\% pop., 10\% land, no tax), 2nd (nobility, 2\% pop., 20\% land, no tax), 3rd (bourgeoisie + peasants, 97\% pop., 70\% land, all taxes + tithes).
    \item \textbf{Causes of French Revolution:} Enlightenment ideas + Madame Deficit + Seven Years' War + famine; Estates General $\rightarrow$ Tennis Court Oath $\rightarrow$ National Assembly $\rightarrow$ Declaration of the Rights of Man $\rightarrow$ secular state/republic.
    \item \textbf{French Revolution gave modernity:} Democracy, the seed of communism, nationalism, citizenship.
    \item \textbf{Modernity blueprint:} Humanism (individualism + secularism), rationality, empiricism, science, reason, democracy, free speech, gender equality, nationalism, citizenship, republic, nuclear family, industrial economy/urbanisation/class.
    \item \textbf{Industrial Revolution (from 1750, ongoing):} Agrarian/feudal $\rightarrow$ industrial/capitalist; cottage industry $\rightarrow$ heavy industry; simple $\rightarrow$ complex division of labour; water/wind $\rightarrow$ coal/steam; subsistence $\rightarrow$ profit economy; unified $\rightarrow$ separated home/work; joint $\rightarrow$ nuclear family.
    \item \textbf{Social evils of industrialisation:} Slums, class inequality, overcrowding, pollution, worker exploitation, degradation of (craft) work---Marx responded with the \textit{Communist Manifesto}.
    \item \textbf{Counter-Enlightenment/conservatives (de Bonald, de Maistre = extreme; Comte, Durkheim, Saint-Simon):} Sought restoration via reform, alarmed by post-Revolution disorder (Reign of Terror); Durkheim specifically opposed anti-Semitism, church-state separation, and urbanisation's effect on family.
    \item \textbf{Radicals (Marx):} Sought revolution, not reform, believing reform merely delays real change.
    \item \textbf{Sociology emerged:} To find social laws (positivism, Comte) that could restore social order and to study modernity's social evils; structural functionalism (reform) and Marxism (revolution) are both, in this sense, positivist philosophies.
\end{itemize}
\end{QuickRevision}